\section{Introducción: LLM como herramienta para el descubrimiento científico y matemático}

Esta clase está impartida por el profesor de UT Austin Swarat Chaudhuri, y se centra en presentar las aplicaciones de los LLM en el descubrimiento matemático y el descubrimiento científico, con énfasis particular en cómo la capacidad de abstracción de los LLM acelera el proceso de descubrimiento.

\begin{knowledgebox}{Proceso del descubrimiento matemático y del descubrimiento científico}
\textbf{Descubrimiento matemático}: modelado $\to$ conjetura $\to$ razonamiento riguroso (demostración) $\to$ contraejemplo/nueva conjetura $\to$ ciclo iterativo.\\
\textbf{Descubrimiento científico}: construcción de teoría $\to$ planteamiento de hipótesis $\to$ diseño experimental $\to$ verificación experimental $\to$ nuevos datos $\to$ nueva teoría.\\
AI for Math busca automatizar el proceso de demostración; AI for Science busca automatizar la generación de hipótesis y el diseño experimental.
\end{knowledgebox}

\subsection{Las tres ventajas principales del LLM Agent en el descubrimiento}

\begin{enumerate}
    \item \textbf{Orientación de la búsqueda}: los LLM contienen internamente una gran cantidad de conocimiento previo, que puede orientar la dirección de la búsqueda en un espacio de hipótesis enorme
    \item \textbf{Aprendizaje por experiencia}: el LLM Agent puede interactuar con el entorno y aprender de la experiencia (incorporándola al prompt o ajustando los pesos mediante RL)
    \item \textbf{Capacidad de abstracción}: los LLM pueden descubrir y construir nuevos conceptos y nuevas herramientas—esta es una dirección totalmente nueva pero con gran potencial
\end{enumerate}
